\subsection*{Актуальность темы}

Аналитические СУБД и системы поддержки принятия решений традиционно оцениваются
с помощью воспроизводимых бенчмарков, задающих как структуру данных, так и
характер рабочей нагрузки. Одним из наиболее распространенных стандартов такого
типа является TPC-H, определяющий фиксированный набор таблиц, правила
генерации данных и набор аналитических запросов для сравнения систем между
собой~\cite{tpch-spec}. Значимость воспроизводимых генераторов данных в этой
области связана не только с корректностью самих измерений, но и с возможностью
масштабирования экспериментов на многопроцессорных и распределенных
платформах~\cite{dewitt-gray}.

Классические средства подготовки набора TPC-H исторически ориентированы на
вывод строковых файлов формата tbl, которые затем требуется преобразовывать в
целевые форматы хранения или загружать в СУБД отдельным конвейером
постобработки~\cite{dbgen-mirror}. Такой подход плохо согласуется с современной
инфраструктурой аналитической обработки данных, в которой широко используются
внутрипамятные колонoчные представления Apache Arrow~\cite{apache-arrow},
колонoчные форматы Parquet~\cite{apache-parquet},
ORC~\cite{apache-orc}, Lance~\cite{lance-format} и
Vortex~\cite{vortex-format}, а также современные lakehouse-системы и
их табличные слои, например Apache Iceberg~\cite{apache-iceberg}.

В этих условиях возникает потребность в генераторе, который сочетает
детерминированность эталонного TPC-H~\cite{tpch-spec}, параллельную генерацию,
прямую запись в современные колонoчные форматы и возможность выгрузки данных не
только на локальный диск, но и в S3-совместимые объектные хранилища.
Дополнительное значение имеет отделение модуля генерации от слоя сериализации,
поскольку такая архитектура позволяет использовать генератор не только как
автономную утилиту, но и как встраиваемый источник батчей
Apache Arrow~\cite{apache-arrow}.

\subsection*{Цель и задачи работы}

Целью настоящей выпускной квалификационной работы является разработка генератора
данных Schengen~\cite{schengen-repo} для аналитического бенчмарка
TPC-H~\cite{tpch-spec},
обеспечивающего прямое формирование данных в современных форматах хранения и
поддерживающего сценарии встроенной генерации в оперативной памяти.

Для достижения поставленной цели необходимо решить следующие задачи:
\begin{enumerate}
    \item проанализировать предметную область аналитических бенчмарков и
    существующие средства генерации данных TPC-H~\cite{tpch-spec};
    \item сформулировать требования к архитектуре генератора, ориентированного
    на современные аналитические форматы и сценарии хранения;
    \item разработать модуль детерминированной генерации таблиц TPC-H и слой
    сериализации в поддерживаемые форматы;
    \item обеспечить возможность прямой интеграции генератора с аналитическими
    системами без промежуточной записи данных на диск;
    \item проверить работоспособность выбранных решений на практике, в том
    числе на примере интеграции с DuckDB~\cite{schengen-duckdb-repo}.
\end{enumerate}

\subsection*{Объект и предмет исследования}

Объектом исследования являются программные средства подготовки данных для
аналитических бенчмарков.

Предметом исследования являются методы детерминированной генерации таблиц
TPC-H~\cite{tpch-spec}, их представления в колонoчной форме и последующей сериализации в
современные форматы хранения.

\subsection*{Практическая значимость работы}

Практическая значимость работы состоит в том, что разработанный генератор
позволяет сократить число промежуточных этапов при подготовке наборов данных
для аналитических экспериментов, упростить воспроизведение тестов и расширить
сценарии применения генератора за счет работы с локальными и удаленными
хранилищами, а также за счет встраивания в OLAP-системы и современные
lakehouse-системы.

\subsection*{Структура работы}

Структура работы включает введение, главы, посвященные анализу предметной
области, проектированию и реализации системы, экспериментальной оценке,
заключение и список использованных источников.
